\documentclass[12pt,a4paper]{article}

\usepackage[T2A]{fontenc}
\usepackage[utf8]{inputenc}
\usepackage[russian]{babel}
\usepackage{amsmath,amssymb,amsthm}
\usepackage{enumitem}
\usepackage{array}
\usepackage{booktabs}
\usepackage{tikz}
\usepackage{geometry}
\usepackage{hyperref}

\geometry{left=25mm,right=15mm,top=20mm,bottom=20mm}

\title{Билет №58 \\ \small Процесс разработки ПО и организация работы команды. Водопадная модель, Vмодель, инкрементальная, спиральная модели. (СпБГУ) \\ Жизненный цикл программного продукта. Отличия альфа, бета и продуктового релиза. (СпБПУ)}
\author{Сериков Владислав Александрович}
\date{16 мая 2026}

\begin{document}

\maketitle
\tableofcontents
\newpage

\section{Базовые понятия}

\subsection{Программное обеспечение и программный продукт}

\textbf{Программное обеспечение} --- это совокупность программ, данных, документации, конфигураций и процедур, предназначенных для решения задач пользователя или организации.

\textbf{Программный продукт} --- это программное обеспечение, рассматриваемое как поставляемый и сопровождаемый результат инженерной деятельности. Программный продукт обычно включает:
\begin{itemize}
    \item исполняемый код или исходный код;
    \item пользовательскую и техническую документацию;
    \item тесты и тестовые данные;
    \item средства установки, настройки и обновления;
    \item лицензии, инструкции эксплуатации и сопровождения;
    \item инфраструктурные артефакты: CI/CD-конфигурации, контейнеры, скрипты миграции и т.д.
\end{itemize}

\textbf{Разработка программного обеспечения} --- это инженерный процесс создания, изменения, проверки, внедрения и сопровождения программного продукта.

\subsection{Процесс разработки ПО}

\textbf{Процесс разработки ПО} --- это упорядоченная совокупность действий, ролей, артефактов, методов и правил, с помощью которых команда создает и сопровождает программный продукт.

Процесс разработки отвечает на вопросы:
\begin{itemize}
    \item какие работы нужно выполнить;
    \item кто за них отвечает;
    \item в каком порядке выполняются работы;
    \item какие результаты должны быть получены;
    \item как проверяется качество;
    \item как управляются изменения, риски и сроки.
\end{itemize}

Типичные группы процессов:
\begin{enumerate}
    \item \textbf{Инженерные процессы:}
    \begin{itemize}
        \item анализ предметной области;
        \item сбор и спецификация требований;
        \item проектирование архитектуры;
        \item детальное проектирование;
        \item программирование;
        \item тестирование;
        \item интеграция;
        \item выпуск;
        \item сопровождение.
    \end{itemize}

    \item \textbf{Управленческие процессы:}
    \begin{itemize}
        \item планирование;
        \item оценка сроков и стоимости;
        \item управление рисками;
        \item управление качеством;
        \item управление конфигурациями;
        \item управление изменениями;
        \item контроль исполнения.
    \end{itemize}

    \item \textbf{Организационные процессы:}
    \begin{itemize}
        \item подбор и обучение команды;
        \item определение стандартов разработки;
        \item настройка инструментов;
        \item совершенствование процесса;
        \item накопление и переиспользование знаний.
    \end{itemize}
\end{enumerate}

\subsection{Артефакты разработки}

\textbf{Артефакт} --- любой значимый результат процесса разработки.

Примеры артефактов:
\begin{itemize}
    \item видение продукта;
    \item требования;
    \item спецификация интерфейсов;
    \item архитектурное описание;
    \item диаграммы UML;
    \item прототипы пользовательского интерфейса;
    \item исходный код;
    \item тест-кейсы;
    \item отчеты о тестировании;
    \item документация пользователя;
    \item релизные заметки;
    \item журнал изменений.
\end{itemize}

\subsection{Качество программного продукта}

Качество ПО обычно рассматривается не только как отсутствие ошибок, но и как соответствие требованиям и ожиданиям заинтересованных сторон.

К важным характеристикам качества относятся:
\begin{itemize}
    \item функциональная пригодность;
    \item надежность;
    \item производительность;
    \item удобство использования;
    \item безопасность;
    \item сопровождаемость;
    \item переносимость;
    \item совместимость.
\end{itemize}

\section{Организация работы команды разработки}

\subsection{Команда разработки}

\textbf{Команда разработки ПО} --- это группа специалистов, совместно выполняющих работы по созданию, проверке, поставке и сопровождению программного продукта.

В зависимости от размера проекта один человек может совмещать несколько ролей.

\subsection{Типовые роли в команде}

\begin{center}
\begin{tabular}{p{0.27\textwidth} p{0.62\textwidth}}
\toprule
\textbf{Роль} & \textbf{Основные обязанности} \\
\midrule
Заказчик & Формулирует бизнес-цели, принимает ключевые результаты, финансирует работу. \\
Пользователь & Использует продукт, дает обратную связь, участвует в проверке удобства и полезности. \\
Product Owner & Управляет видением продукта, приоритизирует требования, ведет backlog. \\
Project Manager & Планирует сроки, ресурсы, бюджет, координирует выполнение работ. \\
Business Analyst / System Analyst & Выявляет, уточняет и документирует требования. \\
Architect & Определяет архитектуру, ключевые технические решения и ограничения. \\
Developer & Реализует функциональность, исправляет дефекты, участвует в code review. \\
QA Engineer / Tester & Планирует и выполняет тестирование, оформляет дефекты, контролирует качество. \\
DevOps / SRE & Настраивает сборку, поставку, мониторинг, инфраструктуру и эксплуатационные процессы. \\
UX/UI Designer & Проектирует пользовательские сценарии и интерфейсы. \\
Technical Writer & Готовит пользовательскую и техническую документацию. \\
\bottomrule
\end{tabular}
\end{center}

\subsection{Коммуникация в команде}

Ключевая проблема разработки ПО состоит в том, что программный продукт является сложной интеллектуальной системой, а требования к нему часто меняются. Поэтому эффективная коммуникация является частью инженерного процесса.

Основные формы коммуникации:
\begin{itemize}
    \item регулярные совещания;
    \item планирование итераций;
    \item ежедневные синхронизации;
    \item демонстрации результата;
    \item ретроспективы;
    \item code review;
    \item обсуждение архитектурных решений;
    \item ведение документации;
    \item работа с системой управления задачами.
\end{itemize}

\subsection{Управление задачами}

Задачи обычно фиксируются в системе управления проектом. Каждая задача должна иметь:
\begin{itemize}
    \item уникальный идентификатор;
    \item краткое название;
    \item описание;
    \item исполнителя;
    \item приоритет;
    \item статус;
    \item критерии готовности;
    \item связи с другими задачами;
    \item оценку трудоемкости.
\end{itemize}

Типичные статусы задачи:
\[
\text{Backlog} \rightarrow \text{To Do} \rightarrow \text{In Progress} \rightarrow \text{Code Review} \rightarrow \text{Testing} \rightarrow \text{Done}.
\]

\subsection{Definition of Done}

\textbf{Definition of Done} --- это согласованный командой набор условий, при выполнении которых задача считается завершенной.

Пример:
\begin{itemize}
    \item код написан;
    \item код проходит автоматические тесты;
    \item выполнено code review;
    \item нет критических дефектов;
    \item документация обновлена;
    \item функциональность развернута на тестовом стенде;
    \item заказчик или Product Owner принял результат.
\end{itemize}

\subsection{Управление конфигурациями}

\textbf{Управление конфигурациями} --- это процесс идентификации, контроля и учета изменений в артефактах проекта.

Основные элементы:
\begin{itemize}
    \item система контроля версий;
    \item правила ветвления;
    \item управление зависимостями;
    \item сборочные сценарии;
    \item хранение артефактов сборки;
    \item маркировка версий;
    \item управление окружениями.
\end{itemize}

Пример простой стратегии ветвления:
\[
\texttt{main} \leftarrow \texttt{develop} \leftarrow \texttt{feature/*}.
\]

Минимальный сценарий:
\begin{enumerate}
    \item Разработчик создает ветку \texttt{feature/login}.
    \item Реализует функциональность авторизации.
    \item Отправляет изменения в удаленный репозиторий.
    \item Создает merge request или pull request.
    \item Другие разработчики выполняют code review.
    \item CI запускает тесты.
    \item После успешной проверки ветка вливается в \texttt{develop}.
\end{enumerate}

\section{Модели жизненного цикла разработки ПО}

\subsection{Понятие модели жизненного цикла}

\textbf{Модель жизненного цикла ПО} --- это абстрактное описание структуры процесса разработки, определяющее основные стадии, порядок их выполнения, правила перехода между ними и типичные результаты каждой стадии.

Модель не является конкретной методикой управления проектом. Она задает общую схему организации работ.

Классические модели:
\begin{itemize}
    \item водопадная модель;
    \item V-модель;
    \item инкрементальная модель;
    \item итерационная модель;
    \item спиральная модель;
    \item гибкие модели разработки.
\end{itemize}

\section{Водопадная модель}

\subsection{Идея модели}

\textbf{Водопадная модель} --- это последовательная модель разработки, в которой стадии выполняются одна за другой, а переход к следующей стадии предполагает завершение предыдущей.

Типичная последовательность:
\[
\text{Требования}
\rightarrow
\text{Проектирование}
\rightarrow
\text{Реализация}
\rightarrow
\text{Тестирование}
\rightarrow
\text{Внедрение}
\rightarrow
\text{Сопровождение}.
\]

\subsection{Схема водопадной модели}

\begin{center}
\begin{tikzpicture}[node distance=1.3cm]
\tikzstyle{block} = [rectangle, draw, rounded corners, minimum width=6cm, minimum height=0.8cm, align=center]

\node[block] (req) {Анализ и спецификация требований};
\node[block, below of=req] (design) {Проектирование};
\node[block, below of=design] (impl) {Реализация};
\node[block, below of=impl] (test) {Тестирование};
\node[block, below of=test] (deploy) {Внедрение};
\node[block, below of=deploy] (support) {Сопровождение};

\draw[->] (req) -- (design);
\draw[->] (design) -- (impl);
\draw[->] (impl) -- (test);
\draw[->] (test) -- (deploy);
\draw[->] (deploy) -- (support);
\end{tikzpicture}
\end{center}

\subsection{Этапы водопадной модели}

\begin{enumerate}
    \item \textbf{Анализ требований.}
    \begin{itemize}
        \item Выявляются потребности заказчика.
        \item Формируется спецификация требований.
        \item Требования согласуются с заинтересованными сторонами.
    \end{itemize}

    \item \textbf{Проектирование.}
    \begin{itemize}
        \item Определяется архитектура системы.
        \item Проектируются модули, интерфейсы, базы данных.
        \item Готовится технический проект.
    \end{itemize}

    \item \textbf{Реализация.}
    \begin{itemize}
        \item Разработчики пишут код.
        \item Создаются модули системы.
        \item Выполняются первичные проверки.
    \end{itemize}

    \item \textbf{Тестирование.}
    \begin{itemize}
        \item Проверяется соответствие продукта требованиям.
        \item Выявляются дефекты.
        \item Исправления возвращаются разработчикам.
    \end{itemize}

    \item \textbf{Внедрение.}
    \begin{itemize}
        \item Система устанавливается в целевой среде.
        \item Пользователи получают доступ к продукту.
        \item Выполняется приемка.
    \end{itemize}

    \item \textbf{Сопровождение.}
    \begin{itemize}
        \item Исправляются ошибки эксплуатации.
        \item Выпускаются обновления.
        \item Обеспечивается поддержка пользователей.
    \end{itemize}
\end{enumerate}

\subsection{Алгоритм применения водопадной модели}

\begin{enumerate}
    \item Зафиксировать цели проекта.
    \item Собрать и согласовать требования.
    \item Подготовить спецификацию требований.
    \item Выполнить архитектурное и детальное проектирование.
    \item Реализовать систему согласно проекту.
    \item Провести тестирование.
    \item Исправить обнаруженные дефекты.
    \item Передать продукт заказчику.
    \item Организовать сопровождение.
\end{enumerate}

\subsection{Минимальный пример}

Разрабатывается простая библиотечная система.

\begin{enumerate}
    \item На этапе требований фиксируется, что система должна хранить книги, читателей и выдачи.
    \item На этапе проектирования создается схема базы данных:
    \[
    \texttt{Book(id, title, author)}, \quad
    \texttt{Reader(id, name)}, \quad
    \texttt{Loan(id, book\_id, reader\_id)}.
    \]
    \item На этапе реализации пишутся модули добавления книг и регистрации выдач.
    \item На этапе тестирования проверяется, что нельзя выдать одну книгу двум читателям одновременно.
    \item После внедрения библиотекари начинают пользоваться системой.
\end{enumerate}

\subsection{Достоинства}

\begin{itemize}
    \item Простота понимания и управления.
    \item Четкая документация.
    \item Хорошо подходит для проектов со стабильными требованиями.
    \item Удобна для контрактной разработки, где результат заранее определен.
    \item Позволяет формально контролировать прохождение стадий.
\end{itemize}

\subsection{Недостатки}

\begin{itemize}
    \item Позднее получение работающего продукта.
    \item Сложность внесения изменений после утверждения требований.
    \item Высокий риск обнаружить ошибки требований только на поздних стадиях.
    \item Слабая адаптация к неопределенности.
    \item Пользователь долго не видит промежуточного результата.
\end{itemize}

\subsection{Область применимости}

Водопадная модель уместна, если:
\begin{itemize}
    \item требования хорошо известны заранее;
    \item предметная область стабильна;
    \item проект относительно формален;
    \item критична документация;
    \item изменения дороги или нежелательны;
    \item продукт создается в регулируемой среде.
\end{itemize}

\section{V-модель}

\subsection{Идея V-модели}

\textbf{V-модель} --- это развитие водопадной модели, в котором каждой стадии проектирования и разработки соответствует стадия верификации и валидации.

Левая ветвь ``V'' описывает уточнение требований и проектирование, нижняя точка --- реализацию, правая ветвь --- тестирование и приемку.

\subsection{Верификация и валидация}

\textbf{Верификация} отвечает на вопрос:
\[
\text{``Правильно ли мы строим продукт?''}
\]
То есть проверяется соответствие промежуточных результатов спецификациям.

\textbf{Валидация} отвечает на вопрос:
\[
\text{``Правильный ли продукт мы строим?''}
\]
То есть проверяется соответствие продукта реальным потребностям пользователя.

\subsection{Схема V-модели}

\begin{center}
\begin{tikzpicture}[node distance=1.3cm]
\tikzstyle{block} = [rectangle, draw, rounded corners, minimum width=4.2cm, minimum height=0.8cm, align=center]

\node[block] (req) at (0,0) {Требования пользователя};
\node[block] (sys) at (1,-1.3) {Системные требования};
\node[block] (arch) at (2,-2.6) {Архитектура};
\node[block] (design) at (3,-3.9) {Детальное проектирование};
\node[block] (code) at (4,-5.2) {Кодирование};

\node[block] (unit) at (5,-3.9) {Модульное тестирование};
\node[block] (int) at (6,-2.6) {Интеграционное тестирование};
\node[block] (sysTest) at (7,-1.3) {Системное тестирование};
\node[block] (accept) at (8,0) {Приемочное тестирование};

\draw[->] (req) -- (sys);
\draw[->] (sys) -- (arch);
\draw[->] (arch) -- (design);
\draw[->] (design) -- (code);
\draw[->] (code) -- (unit);
\draw[->] (unit) -- (int);
\draw[->] (int) -- (sysTest);
\draw[->] (sysTest) -- (accept);

\draw[dashed, <->] (req) -- (accept);
\draw[dashed, <->] (sys) -- (sysTest);
\draw[dashed, <->] (arch) -- (int);
\draw[dashed, <->] (design) -- (unit);
\end{tikzpicture}
\end{center}

\subsection{Соответствие стадий разработки и тестирования}

\begin{center}
\begin{tabular}{p{0.38\textwidth} p{0.45\textwidth}}
\toprule
\textbf{Стадия слева} & \textbf{Соответствующая проверка справа} \\
\midrule
Требования пользователя & Приемочное тестирование \\
Системные требования & Системное тестирование \\
Архитектура & Интеграционное тестирование \\
Детальное проектирование & Модульное тестирование \\
Кодирование & Выполнение тестов \\
\bottomrule
\end{tabular}
\end{center}

\subsection{Алгоритм применения V-модели}

\begin{enumerate}
    \item Описать пользовательские требования.
    \item На их основе заранее определить критерии приемки.
    \item Сформулировать системные требования.
    \item Спроектировать системные тесты.
    \item Разработать архитектуру.
    \item Спроектировать интеграционные тесты.
    \item Выполнить детальное проектирование модулей.
    \item Подготовить модульные тесты.
    \item Реализовать код.
    \item Выполнить модульное тестирование.
    \item Выполнить интеграционное тестирование.
    \item Выполнить системное тестирование.
    \item Выполнить приемочное тестирование.
\end{enumerate}

\subsection{Минимальный пример}

Пусть разрабатывается система интернет-заказа еды.

\begin{itemize}
    \item Пользовательское требование: клиент должен иметь возможность оформить заказ.
    \item Приемочный тест: пользователь выбирает блюда, вводит адрес, оплачивает заказ и получает подтверждение.
    \item Системное требование: система должна сохранять заказ в базе данных и передавать его ресторану.
    \item Системный тест: после оформления заказа запись появляется в таблице заказов и отправляется уведомление ресторану.
    \item Архитектура: отдельные сервисы \texttt{OrderService}, \texttt{PaymentService}, \texttt{NotificationService}.
    \item Интеграционный тест: проверяется совместная работа сервисов заказа и оплаты.
    \item Детальное проектирование: метод \texttt{createOrder()} проверяет корзину, адрес и статус оплаты.
    \item Модульный тест: метод возвращает ошибку при пустой корзине.
\end{itemize}

\subsection{Достоинства}

\begin{itemize}
    \item Явная связь между требованиями и тестами.
    \item Раннее планирование проверки качества.
    \item Хорошая управляемость в проектах с высокими требованиями к надежности.
    \item Подходит для критических систем.
    \item Упрощает трассировку требований.
\end{itemize}

\subsection{Недостатки}

\begin{itemize}
    \item Сохраняет жесткость водопадного подхода.
    \item Требует достаточно стабильных требований.
    \item Изменения могут быть дорогими.
    \item Значительная стоимость документации и тестового планирования.
\end{itemize}

\subsection{Область применимости}

V-модель часто применяется там, где важны:
\begin{itemize}
    \item надежность;
    \item безопасность;
    \item формальная приемка;
    \item соответствие стандартам;
    \item трассируемость требований.
\end{itemize}

Примеры областей:
\begin{itemize}
    \item медицинское ПО;
    \item авиационные и космические системы;
    \item промышленная автоматизация;
    \item банковские системы;
    \item встроенное ПО.
\end{itemize}

\section{Инкрементальная модель}

\subsection{Идея инкрементальной модели}

\textbf{Инкрементальная модель} --- это модель разработки, при которой продукт создается и поставляется частями, называемыми инкрементами.

\textbf{Инкремент} --- это завершенное добавление функциональности к продукту, которое имеет ценность для пользователя или заказчика.

Каждый инкремент может проходить собственные этапы:
\[
\text{Требования}
\rightarrow
\text{Проектирование}
\rightarrow
\text{Реализация}
\rightarrow
\text{Тестирование}
\rightarrow
\text{Поставка}.
\]

\subsection{Схема инкрементальной модели}

\begin{center}
\begin{tikzpicture}
\tikzstyle{block} = [rectangle, draw, rounded corners, minimum width=3.5cm, minimum height=0.8cm, align=center]

\node[block] (i1) at (0,0) {Инкремент 1\\Базовая версия};
\node[block] (i2) at (4,0) {Инкремент 2\\Новая функция};
\node[block] (i3) at (8,0) {Инкремент 3\\Расширение};

\draw[->] (i1) -- (i2);
\draw[->] (i2) -- (i3);

\node at (0,-1.2) {Продукт $P_1$};
\node at (4,-1.2) {Продукт $P_2 = P_1 + \Delta_2$};
\node at (8,-1.2) {Продукт $P_3 = P_2 + \Delta_3$};
\end{tikzpicture}
\end{center}

\subsection{Формальное представление}

Пусть $P_0$ --- начальное состояние продукта, а $\Delta_i$ --- функциональное приращение на шаге $i$. Тогда после $n$ инкрементов продукт можно представить как:
\[
P_n = P_0 + \sum_{i=1}^{n} \Delta_i.
\]

Здесь знак ``$+$'' означает не арифметическое сложение, а добавление функциональности, исправлений и улучшений.

\subsection{Алгоритм инкрементальной разработки}

\begin{enumerate}
    \item Сформировать общее видение продукта.
    \item Разделить требования на группы по приоритетам.
    \item Выбрать минимальный полезный набор функций для первого инкремента.
    \item Реализовать первый инкремент.
    \item Протестировать его.
    \item Передать пользователям или заказчику.
    \item Получить обратную связь.
    \item Уточнить требования к следующим инкрементам.
    \item Повторять шаги 4--8 до достижения целей продукта.
\end{enumerate}

\subsection{Минимальный пример}

Разрабатывается интернет-магазин.

\begin{enumerate}
    \item \textbf{Инкремент 1:} каталог товаров.
    \begin{itemize}
        \item Пользователь видит список товаров.
        \item Может открыть карточку товара.
    \end{itemize}

    \item \textbf{Инкремент 2:} корзина.
    \begin{itemize}
        \item Пользователь добавляет товары в корзину.
        \item Может изменить количество товаров.
    \end{itemize}

    \item \textbf{Инкремент 3:} оформление заказа.
    \begin{itemize}
        \item Пользователь вводит адрес.
        \item Создается заказ.
    \end{itemize}

    \item \textbf{Инкремент 4:} онлайн-оплата.
    \begin{itemize}
        \item Интеграция с платежным сервисом.
        \item Подтверждение оплаты.
    \end{itemize}
\end{enumerate}

После первого инкремента продукт еще не является полноценным магазином, но уже дает часть ценности: можно просматривать товары.

\subsection{Достоинства}

\begin{itemize}
    \item Раннее получение работающего продукта.
    \item Возможность приоритизировать наиболее важные функции.
    \item Снижение риска полного провала проекта.
    \item Получение обратной связи от пользователей.
    \item Возможность постепенного финансирования.
\end{itemize}

\subsection{Недостатки}

\begin{itemize}
    \item Требуется хорошее архитектурное планирование.
    \item Возможен рост технического долга.
    \item Нужна интеграция инкрементов.
    \item Сложно разделить некоторые системы на независимые части.
    \item Частые поставки требуют зрелой инфраструктуры тестирования и сборки.
\end{itemize}

\subsection{Область применимости}

Инкрементальная модель подходит, если:
\begin{itemize}
    \item продукт можно разбить на функциональные части;
    \item важно быстро получить первую версию;
    \item требования могут уточняться;
    \item заказчик готов принимать продукт по частям;
    \item команда способна поддерживать непрерывную интеграцию.
\end{itemize}

\section{Итерационная модель}

\subsection{Связь с инкрементальной разработкой}

Итерационная и инкрементальная модели часто используются вместе, но это разные идеи.

\textbf{Итерация} --- повторение цикла разработки с целью уточнения или улучшения продукта.

\textbf{Инкремент} --- добавление новой функциональности.

Итерационная разработка может улучшать уже существующую функциональность, даже если новых функций не добавляется.

\subsection{Пример различия}

\begin{itemize}
    \item Инкремент: добавить оплату банковской картой.
    \item Итерация: улучшить удобство формы оплаты и уменьшить число ошибок ввода.
\end{itemize}

\section{Спиральная модель}

\subsection{Идея спиральной модели}

\textbf{Спиральная модель} --- это модель разработки, в которой процесс представлен как последовательность циклов, а центральным элементом каждого цикла является анализ и снижение рисков.

Каждый виток спирали уточняет продукт и уменьшает неопределенность.

Спиральная модель сочетает:
\begin{itemize}
    \item итерационность;
    \item прототипирование;
    \item планирование;
    \item анализ рисков;
    \item элементы водопадной модели.
\end{itemize}

\subsection{Основные понятия}

\textbf{Риск} --- потенциальное событие или условие, которое может отрицательно повлиять на цели проекта.

Примеры рисков:
\begin{itemize}
    \item требования неполны или противоречивы;
    \item выбранная технология может не выдержать нагрузку;
    \item интеграция с внешней системой может оказаться сложнее ожидаемой;
    \item пользовательский интерфейс может быть неудобным;
    \item команда может не иметь нужной экспертизы.
\end{itemize}

\subsection{Схема спиральной модели}

Упрощенно каждый виток спирали состоит из четырех квадрантов:

\begin{enumerate}
    \item определение целей, альтернатив и ограничений;
    \item анализ и снижение рисков;
    \item разработка и проверка очередной версии;
    \item планирование следующего витка.
\end{enumerate}

\begin{center}
\begin{tikzpicture}[scale=1.1]
\draw[->] (0,0) -- (4.2,0) node[right] {Разработка};
\draw[->] (0,0) -- (0,4.2) node[above] {Планирование и анализ};

\draw[domain=0:720, smooth, variable=\t, samples=200]
plot ({0.0025*\t*cos(\t)}, {0.0025*\t*sin(\t)});

\node at (2.7,3.5) {1. Цели};
\node at (-2.5,2.5) {2. Риски};
\node at (-2.7,-2.2) {3. Разработка};
\node at (2.8,-1.8) {4. План};

\node at (0,-4.0) {Каждый виток дает более зрелую версию продукта};
\end{tikzpicture}
\end{center}

\subsection{Алгоритм одного витка спирали}

\begin{enumerate}
    \item Определить цели текущего витка.
    \item Определить ограничения: сроки, бюджет, технологии, требования безопасности.
    \item Перечислить альтернативные решения.
    \item Выявить риски для каждой альтернативы.
    \item Оценить вероятность и ущерб рисков.
    \item Выбрать способы снижения наиболее значимых рисков.
    \item При необходимости построить прототип.
    \item Реализовать и проверить очередную версию продукта.
    \item Оценить результаты витка.
    \item Принять решение:
    \begin{itemize}
        \item продолжить разработку;
        \item изменить направление;
        \item повторить исследование;
        \item завершить проект.
    \end{itemize}
    \item Спланировать следующий виток.
\end{enumerate}

\subsection{Минимальный пример}

Разрабатывается система видеоконференций.

\textbf{Виток 1. Проверка технической реализуемости.}
\begin{itemize}
    \item Цель: понять, можно ли обеспечить передачу видео с приемлемой задержкой.
    \item Риск: выбранная технология не выдержит нагрузку.
    \item Действие: создается прототип видеосвязи между двумя пользователями.
    \item Результат: задержка приемлемая, можно продолжать.
\end{itemize}

\textbf{Виток 2. Проверка масштабируемости.}
\begin{itemize}
    \item Цель: поддержать 100 участников.
    \item Риск: серверная инфраструктура слишком дорогая.
    \item Действие: нагрузочное тестирование нескольких архитектур.
    \item Результат: выбирается архитектура с медиасервером.
\end{itemize}

\textbf{Виток 3. Проверка пользовательского опыта.}
\begin{itemize}
    \item Цель: сделать интерфейс создания встречи.
    \item Риск: пользователям сложно подключаться.
    \item Действие: UX-прототип и тестирование на группе пользователей.
    \item Результат: упрощается сценарий входа по ссылке.
\end{itemize}

\subsection{Оценка риска}

Простая количественная оценка риска:
\[
R = P \cdot L,
\]
где:
\begin{itemize}
    \item $R$ --- величина риска;
    \item $P$ --- вероятность наступления риска;
    \item $L$ --- величина потерь при наступлении риска.
\end{itemize}

Например:
\[
P = 0.3,\quad L = 1000000,
\]
тогда:
\[
R = 0.3 \cdot 1000000 = 300000.
\]

Это означает, что ожидаемый ущерб от риска составляет $300000$ условных денежных единиц.

\subsection{Достоинства}

\begin{itemize}
    \item Хорошо работает при высокой неопределенности.
    \item Делает управление рисками центральной частью процесса.
    \item Позволяет использовать прототипирование.
    \item Подходит для крупных и сложных проектов.
    \item Позволяет принимать управленческие решения после каждого витка.
\end{itemize}

\subsection{Недостатки}

\begin{itemize}
    \item Сложна в управлении.
    \item Требует квалифицированного анализа рисков.
    \item Может быть дорогой для малых проектов.
    \item Сложно формально определить момент завершения.
    \item Требует зрелой проектной культуры.
\end{itemize}

\subsection{Область применимости}

Спиральная модель особенно полезна, если:
\begin{itemize}
    \item проект большой и дорогой;
    \item есть существенные технические риски;
    \item требования плохо определены;
    \item нужна проверка гипотез;
    \item возможны разные архитектурные решения;
    \item цена ошибки высока.
\end{itemize}

\section{Сравнение моделей разработки}

\begin{center}
\begin{tabular}{p{0.18\textwidth} p{0.22\textwidth} p{0.22\textwidth} p{0.22\textwidth}}
\toprule
\textbf{Модель} & \textbf{Основная идея} & \textbf{Преимущества} & \textbf{Недостатки} \\
\midrule
Водопадная & Последовательное прохождение стадий & Простота, документация, контроль & Жесткость, поздняя обратная связь \\
V-модель & Связь разработки с тестированием & Трассируемость, раннее планирование тестов & Требует стабильных требований \\
Инкрементальная & Поставка продукта частями & Быстрый результат, обратная связь & Нужна хорошая архитектура \\
Спиральная & Итерации с анализом рисков & Работа с неопределенностью и рисками & Сложность и высокая стоимость управления \\
\bottomrule
\end{tabular}
\end{center}

\section{Жизненный цикл программного продукта}

\subsection{Понятие жизненного цикла}

\textbf{Жизненный цикл программного продукта} --- это совокупность стадий, через которые проходит программный продукт от возникновения идеи до прекращения эксплуатации.

Жизненный цикл шире, чем собственно разработка. Он включает:
\begin{itemize}
    \item замысел;
    \item анализ рынка или потребностей;
    \item проектирование;
    \item реализацию;
    \item тестирование;
    \item выпуск;
    \item эксплуатацию;
    \item сопровождение;
    \item развитие;
    \item вывод из эксплуатации.
\end{itemize}

\subsection{Типовые стадии жизненного цикла}

\begin{enumerate}
    \item \textbf{Инициация.}
    \begin{itemize}
        \item Возникает идея продукта.
        \item Определяются цели и заинтересованные стороны.
        \item Оценивается целесообразность разработки.
    \end{itemize}

    \item \textbf{Анализ требований.}
    \begin{itemize}
        \item Выявляются функциональные и нефункциональные требования.
        \item Определяются ограничения.
        \item Формируется описание продукта.
    \end{itemize}

    \item \textbf{Проектирование.}
    \begin{itemize}
        \item Выбирается архитектура.
        \item Определяются компоненты и интерфейсы.
        \item Проектируются данные и пользовательские сценарии.
    \end{itemize}

    \item \textbf{Реализация.}
    \begin{itemize}
        \item Пишется код.
        \item Создаются конфигурации и миграции.
        \item Выполняется интеграция компонентов.
    \end{itemize}

    \item \textbf{Тестирование.}
    \begin{itemize}
        \item Проверяется соответствие требованиям.
        \item Выполняются модульные, интеграционные, системные, приемочные, нагрузочные и другие тесты.
    \end{itemize}

    \item \textbf{Релиз.}
    \begin{itemize}
        \item Продукт поставляется пользователям.
        \item Публикуются релизные заметки.
        \item Выполняется развертывание.
    \end{itemize}

    \item \textbf{Эксплуатация.}
    \begin{itemize}
        \item Пользователи используют продукт.
        \item Команда отслеживает сбои, производительность и обратную связь.
    \end{itemize}

    \item \textbf{Сопровождение.}
    \begin{itemize}
        \item Исправляются ошибки.
        \item Закрываются уязвимости.
        \item Выпускаются обновления.
    \end{itemize}

    \item \textbf{Развитие.}
    \begin{itemize}
        \item Добавляются новые функции.
        \item Улучшается архитектура.
        \item Оптимизируется производительность.
    \end{itemize}

    \item \textbf{Вывод из эксплуатации.}
    \begin{itemize}
        \item Поддержка прекращается.
        \item Пользователи мигрируют на новую систему.
        \item Данные архивируются или переносятся.
    \end{itemize}
\end{enumerate}

\subsection{Схема жизненного цикла продукта}

\begin{center}
\begin{tikzpicture}[node distance=1.2cm]
\tikzstyle{block} = [rectangle, draw, rounded corners, minimum width=5.5cm, minimum height=0.75cm, align=center]

\node[block] (idea) {Идея и инициация};
\node[block, below of=idea] (req) {Требования};
\node[block, below of=req] (design) {Проектирование};
\node[block, below of=design] (impl) {Реализация};
\node[block, below of=impl] (test) {Тестирование};
\node[block, below of=test] (release) {Релиз};
\node[block, below of=release] (ops) {Эксплуатация и сопровождение};
\node[block, below of=ops] (eol) {Вывод из эксплуатации};

\draw[->] (idea) -- (req);
\draw[->] (req) -- (design);
\draw[->] (design) -- (impl);
\draw[->] (impl) -- (test);
\draw[->] (test) -- (release);
\draw[->] (release) -- (ops);
\draw[->] (ops) -- (eol);

\draw[->, dashed] (ops.east) .. controls +(2,0) and +(2,0) .. (req.east);
\end{tikzpicture}
\end{center}

Пунктирная стрелка показывает, что эксплуатация часто порождает новые требования, после чего продукт снова развивается.

\section{Требования в жизненном цикле ПО}

\subsection{Функциональные требования}

\textbf{Функциональные требования} описывают, что система должна делать.

Примеры:
\begin{itemize}
    \item пользователь может зарегистрироваться;
    \item система рассчитывает стоимость доставки;
    \item администратор может заблокировать учетную запись;
    \item отчет экспортируется в PDF.
\end{itemize}

\subsection{Нефункциональные требования}

\textbf{Нефункциональные требования} описывают свойства системы и ограничения на ее работу.

Примеры:
\begin{itemize}
    \item страница должна открываться не дольше 2 секунд;
    \item система должна выдерживать 10000 одновременных пользователей;
    \item данные должны передаваться по защищенному каналу;
    \item доступность сервиса должна составлять 99.9\%;
    \item интерфейс должен поддерживать русский и английский языки.
\end{itemize}

\subsection{Трассируемость требований}

\textbf{Трассируемость требований} --- способность проследить связь требования с проектными решениями, кодом, тестами и релизами.

Пример трассировки:
\[
\text{REQ-12}
\rightarrow
\text{ARCH-03}
\rightarrow
\text{TASK-145}
\rightarrow
\text{COMMIT-8fa2}
\rightarrow
\text{TEST-77}
\rightarrow
\text{RELEASE-1.4}.
\]

Трассируемость особенно важна в V-модели и в критических системах.

\section{Тестирование в жизненном цикле}

\subsection{Уровни тестирования}

\begin{center}
\begin{tabular}{p{0.3\textwidth} p{0.55\textwidth}}
\toprule
\textbf{Уровень} & \textbf{Что проверяется} \\
\midrule
Модульное тестирование & Отдельные функции, классы, модули. \\
Интеграционное тестирование & Взаимодействие модулей и сервисов. \\
Системное тестирование & Система целиком в условиях, близких к реальным. \\
Приемочное тестирование & Соответствие ожиданиям заказчика или пользователя. \\
\bottomrule
\end{tabular}
\end{center}

\subsection{Виды тестирования}

\begin{itemize}
    \item \textbf{Функциональное тестирование} --- проверка функций системы.
    \item \textbf{Регрессионное тестирование} --- проверка, что изменения не сломали ранее работавшую функциональность.
    \item \textbf{Нагрузочное тестирование} --- проверка поведения под нагрузкой.
    \item \textbf{Стресс-тестирование} --- проверка поведения за пределами нормальной нагрузки.
    \item \textbf{Тестирование безопасности} --- поиск уязвимостей.
    \item \textbf{Usability-тестирование} --- проверка удобства использования.
    \item \textbf{Smoke-тестирование} --- быстрая проверка базовой работоспособности сборки.
\end{itemize}

\section{Релизный цикл программного продукта}

\subsection{Понятие релиза}

\textbf{Релиз} --- это опубликованная или переданная пользователям версия программного продукта.

Релиз может быть:
\begin{itemize}
    \item внутренним;
    \item тестовым;
    \item ограниченным;
    \item публичным;
    \item стабильным;
    \item экспериментальным.
\end{itemize}

\subsection{Версионирование}

Часто используется семантическое версионирование:
\[
MAJOR.MINOR.PATCH.
\]

Например:
\[
2.5.1.
\]

Обычно:
\begin{itemize}
    \item $MAJOR$ увеличивается при несовместимых изменениях;
    \item $MINOR$ увеличивается при добавлении функциональности с сохранением совместимости;
    \item $PATCH$ увеличивается при исправлении ошибок.
\end{itemize}

Пример:
\begin{itemize}
    \item \texttt{1.0.0} --- первый стабильный релиз;
    \item \texttt{1.1.0} --- добавлена новая функция;
    \item \texttt{1.1.1} --- исправлена ошибка;
    \item \texttt{2.0.0} --- изменен API несовместимым образом.
\end{itemize}

\section{Альфа-, бета- и продуктовый релиз}

\subsection{Общая последовательность}

Типичная последовательность релизов:
\[
\text{pre-alpha}
\rightarrow
\text{alpha}
\rightarrow
\text{beta}
\rightarrow
\text{release candidate}
\rightarrow
\text{production release}.
\]

Не все продукты проходят все стадии явно. В небольших проектах некоторые стадии могут объединяться.

\subsection{Альфа-релиз}

\textbf{Альфа-релиз} --- ранняя версия продукта, предназначенная преимущественно для внутреннего тестирования или ограниченного тестирования небольшой группой пользователей.

Характерные признаки:
\begin{itemize}
    \item продукт еще может быть функционально неполным;
    \item интерфейс и архитектура могут существенно меняться;
    \item возможны критические ошибки;
    \item производительность может быть нестабильной;
    \item документация может быть неполной;
    \item основная цель --- проверить концепцию, базовую функциональность и ключевые технические решения.
\end{itemize}

\subsubsection{Цели альфа-релиза}

\begin{itemize}
    \item Проверить, что основные идеи продукта реализуемы.
    \item Найти грубые дефекты.
    \item Получить раннюю обратную связь.
    \item Проверить ключевые архитектурные решения.
    \item Оценить удобство базовых пользовательских сценариев.
\end{itemize}

\subsubsection{Минимальный пример}

Разрабатывается мобильное приложение для учета личных расходов.

Альфа-версия содержит:
\begin{itemize}
    \item добавление расхода;
    \item просмотр списка расходов;
    \item простую локальную базу данных.
\end{itemize}

Но в ней еще нет:
\begin{itemize}
    \item синхронизации с облаком;
    \item отчетов;
    \item экспорта данных;
    \item стабильного дизайна.
\end{itemize}

Такую версию можно дать небольшой группе тестировщиков, чтобы проверить, понятен ли основной сценарий.

\subsection{Бета-релиз}

\textbf{Бета-релиз} --- более зрелая версия продукта, обычно функционально близкая к завершенной, предназначенная для внешнего или расширенного тестирования.

Характерные признаки:
\begin{itemize}
    \item основная функциональность уже реализована;
    \item крупные изменения требований нежелательны;
    \item продукт может использоваться реальными пользователями;
    \item возможны ошибки, но продукт должен быть достаточно стабилен;
    \item цель --- выявить дефекты в реальных условиях эксплуатации;
    \item собирается обратная связь от пользователей.
\end{itemize}

\subsubsection{Цели бета-релиза}

\begin{itemize}
    \item Проверить продукт на реальных сценариях.
    \item Найти ошибки совместимости.
    \item Оценить производительность в реальной среде.
    \item Проверить документацию и процесс установки.
    \item Получить пользовательскую обратную связь перед публичным запуском.
\end{itemize}

\subsubsection{Минимальный пример}

Бета-версия приложения для учета расходов уже содержит:
\begin{itemize}
    \item добавление расходов;
    \item категории;
    \item отчеты;
    \item синхронизацию;
    \item авторизацию;
    \item резервное копирование.
\end{itemize}

Ее дают, например, 1000 пользователям. Пользователи обнаруживают, что на некоторых устройствах синхронизация работает медленно, а отчет за месяц строится с ошибкой при смене часового пояса.

\subsection{Release Candidate}

\textbf{Release Candidate} --- кандидат в финальный релиз. Это версия, которая потенциально может стать продуктовым релизом, если не будут найдены критические ошибки.

Особенности:
\begin{itemize}
    \item новая функциональность обычно уже не добавляется;
    \item исправляются только блокирующие и критические дефекты;
    \item выполняются финальные регрессионные проверки;
    \item проверяется сборка, установка, миграция данных и документация.
\end{itemize}

\subsection{Продуктовый релиз}

\textbf{Продуктовый релиз}, или \textbf{production release}, --- версия продукта, предназначенная для полноценного использования целевыми пользователями в рабочей среде.

Характерные признаки:
\begin{itemize}
    \item продукт считается достаточно стабильным;
    \item реализован заявленный объем функциональности;
    \item критические дефекты устранены или известны и имеют обходные пути;
    \item подготовлена документация;
    \item настроены эксплуатация, мониторинг и поддержка;
    \item определены правила обновления и отката;
    \item продукт доступен целевой аудитории.
\end{itemize}

\subsubsection{Цели продуктового релиза}

\begin{itemize}
    \item Передать продукт пользователям.
    \item Начать полноценную эксплуатацию.
    \item Получить бизнес-эффект.
    \item Организовать поддержку и дальнейшее развитие.
\end{itemize}

\subsection{Сравнение альфа-, бета- и продуктового релиза}

\begin{center}
\begin{tabular}{p{0.2\textwidth} p{0.24\textwidth} p{0.24\textwidth} p{0.24\textwidth}}
\toprule
\textbf{Критерий} & \textbf{Альфа} & \textbf{Бета} & \textbf{Продуктовый релиз} \\
\midrule
Готовность функций & Частичная & Почти полная или полная & Заявленная функциональность готова \\
Аудитория & Внутренняя команда или малая группа & Более широкая группа пользователей & Все целевые пользователи \\
Стабильность & Низкая или средняя & Средняя или высокая & Высокая для эксплуатации \\
Цель & Ранняя проверка идей и архитектуры & Проверка в реальных условиях & Полноценное использование \\
Изменения & Возможны существенные изменения & Обычно только уточнения и исправления & Изменения через новые версии \\
Риск ошибок & Высокий & Умеренный & Минимально допустимый \\
Документация & Неполная & Почти готова & Готова для пользователей \\
\bottomrule
\end{tabular}
\end{center}

\section{Сопровождение программного продукта}

\subsection{Виды сопровождения}

После продуктового релиза жизненный цикл не заканчивается. Значительная часть стоимости ПО связана с сопровождением.

Основные виды сопровождения:
\begin{enumerate}
    \item \textbf{Корректирующее сопровождение.}

    Исправление обнаруженных дефектов.

    \item \textbf{Адаптивное сопровождение.}

    Адаптация продукта к изменениям окружения: операционных систем, библиотек, законодательства, аппаратной платформы.

    \item \textbf{Совершенствующее сопровождение.}

    Добавление новых функций и улучшение существующих.

    \item \textbf{Профилактическое сопровождение.}

    Рефакторинг, оптимизация, устранение технического долга, улучшение тестов.
\end{enumerate}

\subsection{Технический долг}

\textbf{Технический долг} --- накопленные в системе проектные или кодовые решения, которые ускорили разработку в краткосрочной перспективе, но усложняют дальнейшее развитие и сопровождение.

Примеры технического долга:
\begin{itemize}
    \item дублирование кода;
    \item отсутствие тестов;
    \item устаревшие зависимости;
    \item неудачная архитектура;
    \item неявные связи между модулями;
    \item недостаточная документация.
\end{itemize}

\section{Вывод из эксплуатации}

\subsection{Причины вывода из эксплуатации}

Программный продукт может быть выведен из эксплуатации по следующим причинам:
\begin{itemize}
    \item продукт устарел технологически;
    \item поддержка стала слишком дорогой;
    \item появилась новая система;
    \item изменились бизнес-процессы;
    \item закончился срок поддержки;
    \item продукт больше не соответствует требованиям безопасности.
\end{itemize}

\subsection{Типовой алгоритм вывода из эксплуатации}

\begin{enumerate}
    \item Объявить о прекращении поддержки.
    \item Определить дату завершения эксплуатации.
    \item Подготовить план миграции пользователей.
    \item Обеспечить экспорт или перенос данных.
    \item Архивировать необходимые данные.
    \item Отключить интеграции.
    \item Остановить сервис.
    \item Удалить или законсервировать инфраструктуру.
    \item Зафиксировать итоговую документацию.
\end{enumerate}

\section{Связь моделей разработки с жизненным циклом продукта}

Модель разработки описывает, как организована работа по созданию и изменению продукта. Жизненный цикл продукта описывает всю историю продукта от идеи до вывода из эксплуатации.

Один и тот же жизненный цикл может реализовываться разными моделями разработки.

Например:
\begin{itemize}
    \item для первой версии банковской системы может применяться V-модель;
    \item для дальнейшего развития интернет-сервиса может применяться инкрементальная модель;
    \item для исследовательского продукта с высокими рисками может применяться спиральная модель;
    \item для небольшого проекта со стабильными требованиями может быть достаточно водопадной модели.
\end{itemize}

\section{Краткие итоговые положения}

\begin{enumerate}
    \item Процесс разработки ПО включает инженерные, управленческие и организационные действия.
    \item Командная работа требует распределения ролей, коммуникации, управления задачами и контроля качества.
    \item Водопадная модель удобна при стабильных требованиях, но плохо переносит изменения.
    \item V-модель связывает каждую стадию разработки с соответствующим уровнем тестирования.
    \item Инкрементальная модель позволяет поставлять продукт частями и получать раннюю обратную связь.
    \item Спиральная модель ориентирована на управление рисками и полезна в сложных неопределенных проектах.
    \item Жизненный цикл программного продукта включает не только разработку, но и эксплуатацию, сопровождение, развитие и вывод из эксплуатации.
    \item Альфа-релиз предназначен для ранней проверки, бета-релиз --- для расширенного тестирования, продуктовый релиз --- для полноценной эксплуатации.
\end{enumerate}

\end{document}
